\section{Conclusion}\label{sec:conclusion}

This paper studies when posterior uncertainty about volatility materially changes the value of a path-dependent commodity option. The framework separates physical-measure learning from risk-neutral valuation and then asks two distinct questions: how much posterior integration changes a price within a specified model, and how large that effect is relative to the volatility specification required by observed option prices.

The synthetic evidence gives a clear answer to the first question. Posterior integration is not generically negligible. Across 175,000 simulated histories and 168 WTI-like contract states, the largest observed absolute PI--PM gap reaches 0.1384. The largest gaps occur with short historical samples, high volatility, long remaining horizons, extreme moneyness, and little realized fixing. In those states, the second-order approximation
\[
\frac12 C^{Q\prime\prime}(\bar\sigma)\operatorname{Var}(\sigma\mid\mathcal D)
\]
tracks the realized integration gap almost one-for-one. The mechanism is therefore simple but economically meaningful: posterior uncertainty matters when dispersion and pricing curvature are simultaneously large.

The observed WTI application lies in a different region of that state space. In the October-2026 baseline, posterior-integrated and posterior-mean plug-in prices are nearly identical. Their pooled RMSEs are 0.5177 and 0.5183 across 310 observations, and 0.1864 and 0.1869 in the six positive-volume contract-date observations. Average absolute PI--PM difference in that strict subset is about 0.00059. Multi-million-path pseudo-random Monte Carlo, randomized Sobol sequences, and Curran conditioning reproduce these sub-cent differences with numerical uncertainty several orders of magnitude smaller, so the result is not a simulation artifact.

The larger empirical discrepancy is volatility specification. The median historical posterior mean is approximately 0.4155, whereas the median full-cross-section APO-implied effective volatility is approximately 0.4577 under the same one-factor pricing map. More importantly, a strict forward-in-time experiment uses only implied-volatility observations from earlier dates and still lowers MAE from 0.3583 to 0.1297--0.1312 and RMSE from 0.4816 to 0.2312--0.2512 over 345 holdout observations. Valuation-date cluster bootstraps show that these improvements are robust in the full and open-interest-filtered panels. The positive-volume subset points in the same direction but remains too small for equally sharp inference.

Heavy-tailed physical-measure returns improve fit modestly without changing this hierarchy. A Student-$t$ specification produces degrees of freedom near 3.8 and a wider volatility posterior, but its pricing improvement is much smaller than the gain from prior-date option-implied volatility. Prior sensitivity is also limited once the historical sample is long, and large-scale simulation-based calibration supports the Gaussian posterior implementation.

The central empirical conclusion is therefore not that Bayesian parameter uncertainty is irrelevant. It is that, for the observed WTI APO states, posterior integration is a second-order refinement relative to the discrepancy between a long-window historical constant-volatility estimate and the effective volatility surface embedded in option prices. The mechanism experiment simultaneously shows when this ordering can reverse: weakly identified volatility combined with a highly curved pricing map can make posterior integration economically material.

Several qualifications remain. The APO-implied volatility surface is extracted from the same option family and is not an independent external $Q$-volatility source. The one-factor constant-volatility model does not reproduce the entire cross-sectional smile, positive-volume observations remain sparse, and the strongest contract-consistent evidence is concentrated in the September and October 2026 expiries. Historical inference also relies on a continuous/front-month return proxy, while discounting and settlement validation can be refined further.

These limitations point toward the next model layer rather than weakening the main comparison. A natural extension would combine an independently observed vanilla CL volatility surface with a richer futures-curve or stochastic-volatility model and then propagate posterior uncertainty in those additional state variables. The present paper establishes the benchmark against which that additional complexity should be judged: Bayesian integration is worth carrying when posterior dispersion meets pricing curvature, but in the current WTI APO data the dominant empirical margin is the volatility state itself.
